%% paper66_information_theory_en.tex
%% Paper 66: Information Theory and Coding Theory
%% Author: Michael Fuhrmann
%% Project: Specialist Mathematics Project, Build 145
%% Date: 2026-03-12

\documentclass{amsart}

%% ---- Packages ----
\usepackage{amsmath, amssymb, amsthm}
\usepackage{mathtools}
\usepackage{hyperref}
\usepackage{enumitem}
\usepackage{geometry}
\geometry{margin=1in}

%% ---- Theorem environments ----
\newtheorem{theorem}{Theorem}[section]
\newtheorem{lemma}[theorem]{Lemma}
\newtheorem{corollary}[theorem]{Corollary}
\newtheorem{proposition}[theorem]{Proposition}
\theoremstyle{definition}
\newtheorem{definition}[theorem]{Definition}
\newtheorem{example}[theorem]{Example}
\newtheorem{remark}[theorem]{Remark}
\theoremstyle{remark}
\newtheorem{conjecture}[theorem]{Conjecture}
\newtheorem{openquestion}[theorem]{Open Question}

%% ---- Operators ----
\DeclareMathOperator{\Ent}{H}
\DeclareMathOperator{\MI}{I}
\DeclareMathOperator{\KL}{D_{KL}}
\newcommand{\E}{\mathbb{E}}
\newcommand{\R}{\mathbb{R}}
\newcommand{\N}{\mathbb{N}}
\newcommand{\Z}{\mathbb{Z}}
\newcommand{\F}{\mathbb{F}}
\newcommand{\SNR}{\mathrm{SNR}}
\newcommand{\wt}{\mathrm{wt}}
\newcommand{\dist}{\mathrm{d}}

%% ---- Title block ----
\title{Information Theory and Coding Theory:\\
Shannon's Mathematical Framework and Its Foundations}

\author{Michael Fuhrmann}
\address{Specialist Mathematics Project, Build 145}
\email{specialist-maths@project.local}
\date{2026-03-12}

%% ---- Abstract (BEFORE \maketitle) ----
\begin{abstract}
We present a rigorous mathematical treatment of information theory and coding theory,
following Claude Shannon's foundational 1948 paper \emph{A Mathematical Theory of
Communication}. Starting from the axiomatic definition of Shannon entropy
$H(X) = -\sum_i p_i \log_2 p_i$, we develop mutual information, the data-processing
inequality, and the two fundamental theorems of Shannon: the source coding theorem
(noiseless channel coding) and the channel coding theorem (noisy channel coding).
We give complete proofs of the Huffman code optimality, the Singleton bound, the
Hamming bound, and the capacity-achieving property of polar codes (Ar\i{}kan 2009).
We further cover linear codes, Reed--Solomon codes, turbo codes, LDPC codes, and
close with algorithmic information theory, where Kolmogorov complexity $K(x)$ is
shown to be uncomputable. All results are marked with their proof status: proven
theorems versus open conjectures.
\end{abstract}

\maketitle

%% ============================================================
\section{Introduction}
%% ============================================================

In 1948, Claude Elwood Shannon published \emph{A Mathematical Theory of
Communication}~\cite{Shannon1948}, a work that simultaneously founded information
theory as a mathematical discipline and established the engineering science of
digital communication. Shannon's insight was to quantify \emph{information} as a
measurable quantity, independent of the semantic content of messages, and to
characterise the fundamental limits of reliable communication over noisy channels.

Before Shannon, communication engineering relied on heuristics and empirical
design. Shannon demonstrated that every communication channel has a precise
capacity $C$ (measured in bits per channel use), and that reliable communication
at any rate $R < C$ is achievable, while rates $R > C$ are not. This result—the
Channel Coding Theorem—is one of the most profound in all of applied mathematics.

\subsection*{Scope of this paper}

This paper covers:
\begin{enumerate}[label=(\arabic*)]
  \item Shannon entropy and its basic properties (proved);
  \item Mutual information and the data-processing inequality (proved);
  \item The source coding theorem and Huffman code optimality (proved);
  \item Channel capacity and the AWGN Shannon--Hartley formula (proved);
  \item The channel coding theorem and its strong converse (proved);
  \item Error-correcting codes: Hamming and Singleton bounds (proved);
  \item Linear codes, Hamming codes, Reed--Solomon codes;
  \item Capacity-achieving codes: turbo, LDPC, polar (Ar\i{}kan, proved);
  \item Algorithmic information theory and Kolmogorov complexity (partly proved,
        partly open).
\end{enumerate}

All logarithms are base 2 unless otherwise stated; the unit of information is
the \emph{bit}.

%% ============================================================
\section{Entropy}
%% ============================================================

\subsection{Definition and Basic Properties}

\begin{definition}[Shannon Entropy]
Let $X$ be a discrete random variable taking values in a finite alphabet
$\mathcal{X}$ with probability mass function $p(x) = \Pr[X = x]$.
The \emph{Shannon entropy} of $X$ is
\begin{equation}\label{eq:entropy}
  H(X) \;=\; -\sum_{x \in \mathcal{X}} p(x)\log_2 p(x),
\end{equation}
where we adopt the convention $0 \log_2 0 = 0$ (justified by continuity).
The unit is \emph{bits}.
\end{definition}

\begin{remark}
Equivalently, $H(X) = \E[-\log_2 p(X)]$, i.e.\ the expected surprisal.
The \emph{surprisal} (self-information) of an event $x$ is $I(x) = -\log_2 p(x)$.
\end{remark}

\begin{example}
A fair coin ($p = 1/2$) has $H = 1$ bit. A biased coin with $p = 0.9$ has
$H = -0.9\log_2 0.9 - 0.1\log_2 0.1 \approx 0.469$ bits.
\end{example}

\begin{theorem}[Non-negativity]\label{thm:nonneg}
$H(X) \geq 0$, with equality if and only if $X$ is deterministic (i.e.\ $p(x^*)=1$
for some $x^*$).
\end{theorem}
\begin{proof}
Since $0 \leq p(x) \leq 1$, we have $-\log_2 p(x) \geq 0$ for each $x$ with
$p(x)>0$. Hence $H(X) \geq 0$. Equality holds iff $-p(x)\log_2 p(x)=0$ for all
$x$, i.e.\ $p(x) \in \{0,1\}$ for all $x$.
\end{proof}

\begin{theorem}[Maximum Entropy for Uniform Distribution]\label{thm:maxent}
Among all distributions on an alphabet $\mathcal{X}$ of size $n$, entropy is
maximised by the uniform distribution $p(x) = 1/n$, giving $H(X) = \log_2 n$.
\end{theorem}
\begin{proof}
By the log-sum inequality (or Jensen's inequality applied to the concave function
$-t\log_2 t$), for any distribution $p$ on $\mathcal{X}$:
\[
  H(X) = -\sum_x p(x)\log_2 p(x) \leq \log_2 n,
\]
with equality iff $p(x) = 1/n$ for all $x$. Concretely:
\[
  \log_2 n - H(X) = \sum_x p(x)\log_2\!\frac{p(x)}{1/n} = \KL(p \,\|\, u) \geq 0,
\]
where $u$ is the uniform distribution and $\KL \geq 0$ by the non-negativity of
KL-divergence (Gibbs' inequality).
\end{proof}

\subsection{Joint and Conditional Entropy}

\begin{definition}
For jointly distributed random variables $X, Y$ with joint pmf $p(x,y)$:
\begin{align}
  H(X,Y) &= -\sum_{x,y} p(x,y)\log_2 p(x,y), \\
  H(X \mid Y) &= -\sum_{x,y} p(x,y)\log_2 p(x \mid y) = H(X,Y) - H(Y).
\end{align}
\end{definition}

\begin{theorem}[Subadditivity of Entropy]\label{thm:subadd}
$H(X,Y) \leq H(X) + H(Y)$, with equality iff $X$ and $Y$ are independent.
\end{theorem}
\begin{proof}
$H(X) + H(Y) - H(X,Y) = \MI(X;Y) \geq 0$
(mutual information is non-negative, proved below as Theorem~\ref{thm:mi-nonneg}).
\end{proof}

\begin{theorem}[Conditioning Reduces Entropy]\label{thm:condred}
$H(X \mid Y) \leq H(X)$, with equality iff $X \perp Y$.
\end{theorem}
\begin{proof}
$H(X) - H(X\mid Y) = \MI(X;Y) \geq 0$.
\end{proof}

%% ============================================================
\section{Mutual Information}
%% ============================================================

\subsection{Definition}

\begin{definition}[Mutual Information]
The \emph{mutual information} between discrete random variables $X$ and $Y$ is
\begin{equation}\label{eq:mi}
  \MI(X;Y) \;=\; H(X) + H(Y) - H(X,Y)
  \;=\; \sum_{x,y} p(x,y)\log_2\frac{p(x,y)}{p(x)p(y)}.
\end{equation}
Equivalently, $\MI(X;Y) = \KL(p_{XY} \,\|\, p_X \otimes p_Y)$.
\end{definition}

\begin{theorem}[Non-negativity of Mutual Information]\label{thm:mi-nonneg}
$\MI(X;Y) \geq 0$, with equality iff $X$ and $Y$ are independent.
\end{theorem}
\begin{proof}
$\MI(X;Y) = \KL(p_{XY}\|p_Xp_Y)$. By Gibbs' inequality (or Jensen applied to
$-\log$):
\[
  \KL(p \| q) = \sum_x p(x)\log_2\frac{p(x)}{q(x)} \geq 0,
\]
with equality iff $p = q$ everywhere. Hence $\MI(X;Y) \geq 0$, with equality iff
$p(x,y) = p(x)p(y)$ for all $x,y$.
\end{proof}

\subsection{Data-Processing Inequality}

\begin{theorem}[Data-Processing Inequality]\label{thm:dpi}
If $X \to Y \to Z$ is a Markov chain (i.e.\ $Z$ is conditionally independent of
$X$ given $Y$), then
\[
  \MI(X;Z) \leq \MI(X;Y).
\]
\end{theorem}
\begin{proof}
By the chain rule of mutual information:
\begin{align*}
  \MI(X;Y,Z) &= \MI(X;Y) + \MI(X;Z\mid Y), \\
  \MI(X;Y,Z) &= \MI(X;Z) + \MI(X;Y\mid Z).
\end{align*}
Since $X \to Y \to Z$ is Markov, $\MI(X;Z \mid Y) = 0$. Hence
$\MI(X;Y) = \MI(X;Z) + \MI(X;Y \mid Z)$.
Because $\MI(X;Y\mid Z) \geq 0$, we obtain $\MI(X;Z) \leq \MI(X;Y)$.
\end{proof}

\begin{remark}
The data-processing inequality expresses that no post-processing of $Y$ can
increase information about $X$. It is fundamental in machine learning (no
deterministic function of features can exceed the information in the features)
and in the study of communication networks.
\end{remark}

%% ============================================================
\section{Shannon's Source Coding Theorem}
%% ============================================================

\subsection{Lossless Compression}

\begin{definition}[Variable-Length Code]
A \emph{prefix-free} (instantaneous) code for $\mathcal{X}$ assigns to each
$x \in \mathcal{X}$ a binary codeword $C(x) \in \{0,1\}^*$ such that no codeword
is a prefix of another. The \emph{expected length} is
$\bar{L} = \sum_x p(x)\,|C(x)|$.
\end{definition}

\begin{theorem}[Source Coding Theorem (Shannon 1948)]\label{thm:source-coding}
For any discrete random variable $X$ with entropy $H(X)$, there exists a
prefix-free code with expected length $\bar{L}$ satisfying
\[
  H(X) \;\leq\; \bar{L} \;<\; H(X) + 1.
\]
Moreover, no prefix-free code can achieve $\bar{L} < H(X)$.
\end{theorem}
\begin{proof}
\textbf{Lower bound.} By Kraft's inequality, any prefix-free code satisfies
$\sum_x 2^{-|C(x)|} \leq 1$. Let $l_x = |C(x)|$. Then:
\begin{align*}
  \bar{L} - H(X) &= \sum_x p(x)(l_x + \log_2 p(x)) \\
  &= \sum_x p(x)\log_2\frac{p(x)}{2^{-l_x}}
   = \KL\!\left(p \,\Big\|\, \{2^{-l_x}\}\right) \geq 0.
\end{align*}

\textbf{Upper bound.} Choose $l_x = \lceil -\log_2 p(x)\rceil$. Then
$l_x < -\log_2 p(x) + 1$, so $\sum_x 2^{-l_x} \leq \sum_x p(x) = 1$ (Kraft
satisfied). The expected length satisfies
\[
  \bar{L} = \sum_x p(x)l_x < \sum_x p(x)(-\log_2 p(x) + 1) = H(X) + 1. \qedhere
\]
\end{proof}

\subsection{Huffman Codes}

\begin{algorithm*}[Huffman's Algorithm (1952)]
Given symbols with probabilities $p_1 \geq p_2 \geq \cdots \geq p_n$:
\begin{enumerate}
  \item Treat each symbol as a leaf node weighted by its probability.
  \item Repeatedly merge the two nodes of lowest weight into a new parent node
        (weight = sum), labelling the two children 0 and 1.
  \item The resulting binary tree defines a prefix-free code.
\end{enumerate}
\end{algorithm*}

\begin{theorem}[Huffman Optimality]\label{thm:huffman}
Among all prefix-free codes, the Huffman code minimises the expected codeword
length $\bar{L}$.
\end{theorem}
\begin{proof}[Proof sketch]
By induction on $n$. Base case $n=2$: both symbols get 1-bit codes; optimal.
Inductive step: Let $C^*$ be any optimal code. Among all symbols with minimum
probability, at least two share the same maximum codeword length (otherwise
we could shorten the longer one, reducing $\bar{L}$). WLOG the two least
probable symbols $x_{n-1}, x_n$ form a sibling pair in $C^*$. Merging them into
a meta-symbol of probability $p_{n-1}+p_n$ gives an $(n-1)$-symbol code $C'$
with $\bar{L}(C^*) = \bar{L}(C') + (p_{n-1}+p_n)$. By the induction hypothesis,
Huffman's greedy merge minimises $\bar{L}(C')$, hence also $\bar{L}(C^*)$.
\end{proof}

%% ============================================================
\section{Channel Capacity}
%% ============================================================

\subsection{Discrete Memoryless Channels}

\begin{definition}[Discrete Memoryless Channel]
A \emph{discrete memoryless channel} (DMC) is specified by input alphabet
$\mathcal{X}$, output alphabet $\mathcal{Y}$, and transition probabilities
$p(y \mid x)$. The channel is \emph{memoryless}: $p(y^n \mid x^n) =
\prod_{i=1}^n p(y_i \mid x_i)$.
\end{definition}

\begin{definition}[Channel Capacity]
The \emph{capacity} of a DMC is
\begin{equation}\label{eq:capacity}
  C \;=\; \max_{p(x)} \MI(X;Y),
\end{equation}
where the maximum is over all input distributions $p(x)$.
\end{definition}

\begin{example}[Binary Symmetric Channel]
A BSC with crossover probability $\varepsilon \in [0,1/2)$ has
$C_{\mathrm{BSC}} = 1 - H_b(\varepsilon)$,
where $H_b(\varepsilon) = -\varepsilon\log_2\varepsilon -
(1-\varepsilon)\log_2(1-\varepsilon)$ is the binary entropy function.
\end{example}

\subsection{The AWGN Channel and Shannon--Hartley Theorem}

\begin{definition}[AWGN Channel]
The \emph{additive white Gaussian noise} (AWGN) channel model is
$Y = X + Z$, where $Z \sim \mathcal{N}(0, N_0/2)$ is independent of the input
$X$, which satisfies a power constraint $\E[X^2] \leq P$.
\end{definition}

\begin{theorem}[Shannon--Hartley Theorem]\label{thm:shannon-hartley}
The capacity of the AWGN channel with signal power $P$ and noise power $N$ is
\begin{equation}\label{eq:awgn}
  C \;=\; \tfrac{1}{2}\log_2\!\left(1 + \tfrac{P}{N}\right)
  \quad [\text{bits per channel use}].
\end{equation}
For a bandwidth $W$ (Hz) and noise power spectral density $N_0/2$, this gives
$C = W \log_2(1 + P/(N_0 W))$ bits per second.
\end{theorem}
\begin{proof}[Proof sketch]
By the entropy-power inequality, the Gaussian distribution maximises differential
entropy $h(Y)$ subject to a power constraint. With $h(Z) = \frac{1}{2}\log_2(2\pi
e N)$ and $h(Y) \leq \frac{1}{2}\log_2(2\pi e(P+N))$, we get $\MI(X;Y) = h(Y) -
h(Z) \leq \frac{1}{2}\log_2(1+P/N)$, achieved by $X \sim \mathcal{N}(0,P)$.
\end{proof}

%% ============================================================
\section{Shannon's Channel Coding Theorem}
%% ============================================================

\begin{theorem}[Channel Coding Theorem (Shannon 1948)]\label{thm:channel-coding}
For a DMC with capacity $C$:
\begin{enumerate}[label=(\alph*)]
  \item \textbf{Achievability.} For every rate $R < C$ and every $\varepsilon > 0$,
        there exist block codes of length $n$ (for sufficiently large $n$) with
        rate $R$ and maximal probability of error $P_e^{(n)} < \varepsilon$.
  \item \textbf{Converse.} If $R > C$, then for every sequence of codes the
        probability of error $P_e^{(n)} \to 1$ as $n \to \infty$.
\end{enumerate}
\end{theorem}
\begin{proof}[Proof sketch (Random Coding Argument)]
\textbf{Achievability.} Generate a codebook $\mathcal{C}$ of $2^{nR}$ codewords
$\mathbf{x}(m)$, $m \in [2^{nR}]$, each drawn i.i.d.\ from the capacity-achieving
input distribution $p^*(x)$. The decoder uses \emph{typical set decoding}:
declare $\hat{m}$ if $(\mathbf{x}(\hat{m}), \mathbf{y})$ is jointly typical and
unique. By the asymptotic equipartition property (AEP), the probability that the
correct codeword is typical is $\geq 1-\varepsilon$ for large $n$. The probability
that any other codeword $m' \neq m$ is jointly typical with $\mathbf{y}$ is
$\leq 2^{-n(I(X;Y)-\varepsilon)}$. By union bound over $2^{nR}$ wrong codewords:
\[
  P_e \leq \varepsilon + 2^{nR} \cdot 2^{-n(C-\varepsilon)} \to 0 \quad \text{if } R < C.
\]

\textbf{Converse.} By Fano's inequality: if $P_e$ is small,
$H(M \mid Y^n) \leq 1 + P_e \cdot nR$. Also $\MI(M; Y^n) \leq nC$ (channel
bound). Since $nR = H(M) = \MI(M;Y^n) + H(M\mid Y^n)$, we get $R \leq C + O(P_e)$,
forcing $P_e \not\to 0$ when $R > C$.
\end{proof}

\begin{theorem}[Strong Converse]\label{thm:strong-converse}
For a DMC with capacity $C$, if $R > C$, then $P_e^{(n)} \to 1$ exponentially fast.
\end{theorem}
\begin{proof}[Proof sketch]
Using the method of types or sphere-packing exponents, the probability of error
satisfies $P_e^{(n)} \geq 1 - 2^{-n(E(R) - o(1))}$ for some positive exponent
$E(R)>0$ when $R>C$, showing $P_e^{(n)} \to 1$.
\end{proof}

%% ============================================================
\section{Error-Correcting Codes}
%% ============================================================

\subsection{Hamming Distance and Basic Definitions}

\begin{definition}[Hamming Distance]
For $\mathbf{u}, \mathbf{v} \in \{0,1\}^n$, the \emph{Hamming distance} is
\[
  \dist_H(\mathbf{u}, \mathbf{v}) = |\{i : u_i \neq v_i\}|.
\]
The \emph{Hamming weight} of $\mathbf{v}$ is $\wt(\mathbf{v}) = \dist_H(\mathbf{v},
\mathbf{0})$.
\end{definition}

\begin{definition}[$[n,k,d]$-Code]
A \emph{binary block code} $\mathcal{C} \subseteq \{0,1\}^n$ with $|\mathcal{C}| =
2^k$ and minimum distance $d = \min_{\mathbf{c} \neq \mathbf{c}'} \dist_H(\mathbf{c},
\mathbf{c}')$ is an \emph{$[n,k,d]$-code}. It can:
\begin{itemize}
  \item detect up to $d-1$ errors,
  \item correct up to $\lfloor (d-1)/2 \rfloor$ errors.
\end{itemize}
The \emph{rate} is $R = k/n$. The \emph{redundancy} is $n - k$ bits.
\end{definition}

\subsection{Singleton Bound}

\begin{theorem}[Singleton Bound]\label{thm:singleton}
Any $[n,k,d]$-code satisfies
\[
  k \;\leq\; n - d + 1.
\]
\end{theorem}
\begin{proof}
Consider the $2^k$ codewords. Project onto the first $k' = n-d+1$ coordinates;
the resulting vectors must all be distinct (otherwise two codewords agree on all
$k'$ positions, and since they have distance $\geq d$, they must differ in at
least $d$ of the remaining $n-k'=d-1$ positions, which is impossible). Hence
$2^k \leq 2^{k'} = 2^{n-d+1}$.
\end{proof}

\begin{definition}[MDS Codes]
Codes achieving $k = n - d + 1$ are called \emph{maximum distance separable}
(MDS). Reed--Solomon codes are MDS.
\end{definition}

\subsection{Hamming (Sphere-Packing) Bound}

\begin{theorem}[Hamming Bound]\label{thm:hamming-bound}
An $[n,k,d]$-code with $d \geq 2t+1$ (correcting $t$ errors) satisfies
\[
  2^k \cdot \sum_{i=0}^{t}\binom{n}{i} \;\leq\; 2^n.
\]
\end{theorem}
\begin{proof}
Each codeword $\mathbf{c}$ has a \emph{Hamming ball} of radius $t$:
$B(\mathbf{c}, t) = \{\mathbf{v} : \dist_H(\mathbf{v},\mathbf{c}) \leq t\}$,
with $|B(\mathbf{c},t)| = \sum_{i=0}^t \binom{n}{i}$. Since $d \geq 2t+1$, these
balls are disjoint, and all are subsets of $\{0,1\}^n$. Counting yields the bound.
\end{proof}

\begin{definition}[Perfect Code]
A code meeting the Hamming bound with equality is called \emph{perfect}.
The Hamming codes (see Section~\ref{sec:linear}) are perfect.
\end{definition}

%% ============================================================
\section{Linear Codes}\label{sec:linear}
%% ============================================================

\subsection{Generator and Parity-Check Matrices}

\begin{definition}[Linear Code]
A \emph{linear $[n,k]$-code} over $\mathbb{F}_2$ is a $k$-dimensional subspace
of $\mathbb{F}_2^n$. It is described by:
\begin{itemize}
  \item A \emph{generator matrix} $G \in \mathbb{F}_2^{k \times n}$: codewords
        are $\mathbf{c} = \mathbf{m} G$ for messages $\mathbf{m} \in \mathbb{F}_2^k$.
  \item A \emph{parity-check matrix} $H \in \mathbb{F}_2^{(n-k) \times n}$ with
        $H G^T = 0$: $\mathbf{c}$ is a codeword iff $H \mathbf{c}^T = \mathbf{0}$.
\end{itemize}
\end{definition}

\begin{definition}[Syndrome]
For a received word $\mathbf{r}$, the \emph{syndrome} is $\mathbf{s} = H
\mathbf{r}^T$. If $\mathbf{s} = \mathbf{0}$, no (detectable) error occurred.
\end{definition}

\subsection{Hamming Codes}

\begin{definition}[Hamming Code $[2^r-1, 2^r-r-1, 3]$]
For $r \geq 2$, the \emph{Hamming code} has parity-check matrix $H$ whose columns
are all $2^r - 1$ nonzero binary vectors of length $r$. Parameters:
$n = 2^r - 1$, $k = 2^r - r - 1$, $d = 3$.
\end{definition}

\begin{proposition}[Hamming Codes are Perfect]
Hamming codes achieve the Hamming bound with equality, hence are perfect 1-error-
correcting codes.
\end{proposition}
\begin{proof}
With $t=1$: $2^k \cdot (1 + n) = 2^{2^r-r-1} \cdot 2^r = 2^{2^r-1} = 2^n$.
\end{proof}

\subsection{Reed--Solomon Codes}

\begin{definition}[Reed--Solomon Code]
Let $q = p^m$ be a prime power and $\alpha_1,\ldots,\alpha_n$ be $n$ distinct
elements of $\mathbb{F}_q$. The \emph{Reed--Solomon code} of dimension $k$ is
\[
  \mathcal{C}_{RS} = \bigl\{(f(\alpha_1),\ldots,f(\alpha_n)) :
  f \in \mathbb{F}_q[x],\; \deg f \leq k-1\bigr\}.
\]
It is an $[n, k, n-k+1]$-code over $\mathbb{F}_q$, hence MDS.
\end{definition}

\begin{remark}
Reed--Solomon codes over $\mathbb{F}_{2^8}$ are used in CD/DVD error correction
and in RAID-6 storage systems. Efficient decoding up to $\lfloor(n-k)/2\rfloor$
errors is possible via the Berlekamp--Massey or Euclidean algorithm.
\end{remark}

%% ============================================================
\section{Capacity-Achieving Codes}
%% ============================================================

\subsection{Turbo Codes}

\begin{definition}[Turbo Code (Berrou--Glavieux--Thitimajshima, 1993)]
A \emph{turbo code} is a parallel concatenation of two recursive systematic
convolutional (RSC) encoders, separated by a pseudo-random interleaver. Decoding
uses iterative \emph{belief propagation} (turbo decoding) between the two
component decoders.
\end{definition}

\begin{remark}
Turbo codes were the first practical codes to approach Shannon capacity for the
AWGN channel. Their capacity-achieving property in the limit of increasing block
length and code rate follows from density evolution analysis, though the original
1993 paper demonstrated near-capacity performance empirically.
\end{remark}

\subsection{LDPC Codes}

\begin{definition}[LDPC Code (Gallager, 1960/1963)]
A \emph{low-density parity-check} (LDPC) code is a linear code whose parity-check
matrix $H$ is sparse (most entries zero). Gallager's original regular LDPC codes
have constant column weight $d_v$ and row weight $d_c$.
\end{definition}

\begin{remark}
LDPC codes decoded by belief propagation on the Tanner graph can achieve rates
arbitrarily close to the Shannon capacity of the AWGN channel. This was shown
by Richardson and Urbanke (2001) via density evolution.
\end{remark}

\subsection{Polar Codes and Ar\i{}kan's Theorem}

\begin{definition}[Polar Code (Ar\i{}kan, 2009)]
Let $W$ be a binary-input memoryless channel with capacity $C(W)$. Polar codes
use the \emph{channel polarisation} phenomenon: applying $n = 2^\ell$ copies of
$W$ through the polarisation transform $G_n = B_n F^{\otimes \ell}$ (where
$F = \begin{pmatrix}1&0\\1&1\end{pmatrix}$ and $B_n$ is a bit-reversal permutation)
creates $n$ synthetic channels $W_n^{(i)}$, $i = 1, \ldots, n$, which polarise
to either perfect ($I(W_n^{(i)}) \to 1$) or useless ($I(W_n^{(i)}) \to 0$) channels.
Information bits are sent only on the ``good'' channels; the rest are frozen.
\end{definition}

\begin{theorem}[Capacity-Achieving Property of Polar Codes (Ar\i{}kan 2009)]
\label{thm:polar}
For any binary-input memoryless symmetric channel $W$ with capacity $C(W)$ and
any rate $R < C(W)$, polar codes of block length $n = 2^\ell$ achieve a block
error probability
\[
  P_e \;\leq\; 2^{-n^\beta}
\]
for any $\beta < 1/2$ and sufficiently large $n$, using encoding and decoding
algorithms of complexity $O(n \log n)$.
\end{theorem}
\begin{proof}[Proof sketch]
Ar\i{}kan's key insight is the polarisation lemma: for $I_n = \frac{1}{n}\sum_i
I(W_n^{(i)})$, we have $I_n = C(W)$ for all $n$ (conservation of capacity). By
the martingale convergence theorem, the fraction of channels with $I(W_n^{(i)}) >
1-\delta$ converges to $C(W)$, and the fraction with $I(W_n^{(i)}) < \delta$
converges to $1-C(W)$. The error probability on each good channel decays as
$2^{-2^{\ell\beta}}$ by the recursive structure of successive cancellation decoding,
from which the claimed bound follows. Details are in \cite{Arikan2009}.
\end{proof}

\begin{remark}
Polar codes are the first provably capacity-achieving codes with efficient
($O(n\log n)$) encoding and decoding. They are adopted in the 5G NR standard
for control channel coding.
\end{remark}

%% ============================================================
\section{Algorithmic Information Theory}
%% ============================================================

\subsection{Kolmogorov Complexity}

\begin{definition}[Kolmogorov Complexity (Solomonoff 1960, Kolmogorov 1965, Chaitin 1969)]
Fix a universal Turing machine $\mathcal{U}$. The \emph{Kolmogorov complexity}
of a finite binary string $x$ is
\[
  K(x) \;=\; \min\bigl\{|\pi| : \mathcal{U}(\pi) = x\bigr\},
\]
the length of the shortest program $\pi$ that causes $\mathcal{U}$ to output $x$.
\end{definition}

\begin{theorem}[Invariance]\label{thm:invariance}
Up to an additive constant depending only on the choice of universal Turing machine,
$K(x)$ is well-defined (independent of $\mathcal{U}$).
\end{theorem}
\begin{proof}
Any universal Turing machine $\mathcal{U}_1$ can simulate $\mathcal{U}_2$ with a
fixed prefix (interpreter) of length $c_{12}$. Hence $K_1(x) \leq K_2(x) +
c_{12}$, and symmetrically $K_2(x) \leq K_1(x) + c_{21}$.
\end{proof}

\begin{theorem}[Uncomputability of Kolmogorov Complexity]\label{thm:kolmogorov-uncomputable}
The function $x \mapsto K(x)$ is not computable.
\end{theorem}
\begin{proof}
Suppose $K$ were computable. Enumerate all binary strings $x_1, x_2, \ldots$ in
lexicographic order, and let $n$ be a large integer. Define the string
$s = $ ``the first string $x$ with $K(x) > n$''. This string $s$ is well-defined
(since $K$ grows without bound) and computable from $n$. The program computing
$s$ has length $O(\log n)$, so $K(s) = O(\log n) \ll n$ for large $n$. But
$K(s) > n$ by construction—a contradiction.
\end{proof}

\subsection{Chaitin's $\Omega$ and Randomness}

\begin{definition}[Chaitin's $\Omega$]
For a prefix-free universal Turing machine, the \emph{halting probability}
\[
  \Omega \;=\; \sum_{\mathcal{U}(\pi)\text{ halts}} 2^{-|\pi|}
\]
is a real number in $(0,1)$ that is \emph{algorithmically random}: its binary
expansion is incompressible in the sense that $K(\Omega \upharpoonright n) \geq n -
O(1)$.
\end{definition}

\begin{openquestion}[Chaitin's $\Omega$]
No explicit closed-form description of $\Omega$ is known. Moreover, no algorithm
can compute more than finitely many bits of $\Omega$; indeed, each bit of $\Omega$
encodes an independent ``oracular'' fact about the halting problem. Whether any
further structural properties of $\Omega$ can be determined remains open.
\end{openquestion}

\begin{theorem}[AIT and Shannon Entropy (Brudno 1983)]
For an ergodic stationary source $\mu$, almost every sequence $x$ satisfies
\[
  \lim_{n\to\infty} \frac{K(x_{1:n})}{n} = H(\mu),
\]
where $H(\mu)$ is the entropy rate of the source.
\end{theorem}
\begin{proof}
This follows from Brudno's theorem, which relates the Kolmogorov complexity
rate of individual sequences to the measure-theoretic entropy rate; see
\cite{CoverThomas2006} for details.
\end{proof}

%% ============================================================
\section{References}
%% ============================================================

\begin{thebibliography}{99}

\bibitem{Shannon1948}
C.~E. Shannon,
\emph{A Mathematical Theory of Communication},
Bell System Technical Journal \textbf{27} (1948), 379--423 and 623--656.

\bibitem{Shannon1959}
C.~E. Shannon,
\emph{Coding Theorems for a Discrete Source with a Fidelity Criterion},
IRE National Convention Record \textbf{4} (1959), 142--163.

\bibitem{CoverThomas2006}
T.~M. Cover and J.~A. Thomas,
\emph{Elements of Information Theory}, 2nd ed.,
Wiley-Interscience, Hoboken, NJ, 2006.

\bibitem{Gallager1963}
R.~G. Gallager,
\emph{Low-Density Parity-Check Codes},
MIT Press, Cambridge, MA, 1963.

\bibitem{Arikan2009}
E.~Ar\i{}kan,
\emph{Channel Polarization: A Method for Constructing Capacity-Achieving Codes
for Symmetric Binary-Input Memoryless Channels},
IEEE Transactions on Information Theory \textbf{55} (2009), no.~7, 3051--3073.

\bibitem{Huffman1952}
D.~A. Huffman,
\emph{A Method for the Construction of Minimum-Redundancy Codes},
Proceedings of the IRE \textbf{40} (1952), no.~9, 1098--1101.

\bibitem{Kolmogorov1965}
A.~N. Kolmogorov,
\emph{Three Approaches to the Quantitative Definition of Information},
Problems of Information Transmission \textbf{1} (1965), no.~1, 1--7.

\bibitem{RichardsonUrbanke2001}
T.~J. Richardson and R.~L. Urbanke,
\emph{The Capacity of Low-Density Parity-Check Codes under Message-Passing
Decoding},
IEEE Transactions on Information Theory \textbf{47} (2001), no.~2, 599--618.

\bibitem{Berrou1993}
C.~Berrou, A.~Glavieux, and P.~Thitimajshima,
\emph{Near Shannon Limit Error-Correcting Coding and Decoding: Turbo-Codes},
Proc.\ IEEE ICC 1993, 1064--1070.

\bibitem{Singleton1964}
R.~C. Singleton,
\emph{Maximum Distance $q$-nary Codes},
IEEE Transactions on Information Theory \textbf{10} (1964), no.~2, 116--118.

\end{thebibliography}

\end{document}
